\section{Introduction}
\label{sec:introduction}

A productive line of cognitive-modelling work asks whether neural
language models (LMs) capture aspects of human language processing: they
predict reading times through surprisal
\citep{levy2008expectation,wilcox2020predictive}, they match syntactic
acceptability patterns \citep{warstadt2020blimp}, and their activations
align with neural recordings \citep{schrimpf2021neural}. An adjacent
question --- whether LMs can model \emph{language acquisition} rather
than steady-state competence --- is only recently being taken up in
earnest \citep{stearns2026artificial}. Companion work shows that a
GPT-2 continually trained on EFCAMDAT learner text produces errors
concentrated in precisely the categories that second-language acquisition
(SLA) research identifies as developmentally vulnerable: verbal
morphology, subject--verb agreement, determiners, and subject pronoun
omission \citep{stearns2026artificial,geertzen2014efcamdat,
goldschneider2001explaining}. The error profile of the ``artificial
learner'' is not uniformly degraded; it is \emph{selectively shifted}
toward the interlanguage attractors predicted by SLA theory.

That result establishes the hypothesis but not the mechanism of control.
The learner GPT-2 is a single model with a fixed, implicit learner
profile shaped by whatever distribution of L1s and proficiencies happens
to dominate EFCAMDAT. At inference, there is no handle that asks it to
predict \emph{as an A2 Spanish-L1 learner} rather than as the corpus
average. Yet for any downstream use of this line of modelling ---
next-word prediction in intelligent tutoring systems, cloze-based
proficiency assessment, generation of calibrated practice material,
diagnostic perplexity mapping \citep{stearns2025nwpaleval} --- the
learner profile is exactly the variable that should be exposed.

We ask: can we add a small architectural knob that (i) conditions
next-word prediction on explicit learner metadata (CEFR proficiency, L1)
and (ii) does so at the mechanism that has the strongest empirical
evidence of being the useful modulation point in softmax attention?

Recent work by the Qwen team \citep{qiu2025gated} systematically
evaluates thirty variants of gating in softmax attention and identifies
a single modification that consistently improves perplexity, downstream
accuracy, training stability, and long-context extrapolation: a
head-specific, elementwise, sigmoid gate applied to the output of the
scaled dot-product attention (SDPA). The gate takes the pre-normalised
hidden state as input and produces a per-dimension filter on the attention
output. Qiu et al.~attribute the gain to two effects: the gate introduces
non-linearity between two otherwise-linear projections, and its
input-dependent sparsity reduces the \emph{attention-sink} phenomenon
\citep{xiao2023efficient}.

Our proposal is a minimal change to this mechanism: extend the gate's
input from the hidden state alone, $X$, to $[X\,;\,e_\text{CEFR}\,;\,e_\text{L1}]$,
the concatenation of the hidden state with learned CEFR and L1
embeddings, broadcast across the sequence. The gating filter thereby
becomes metadata-dependent at every attention layer, while the rest of
the model --- embeddings, feed-forward blocks, layer norms, output head
--- is unchanged. Section~\ref{sec:method} formalises the modification.

\paragraph{Contributions.}
\begin{enumerate}
  \item A metadata-aware extension of Qwen-style SDPA-output gating
    suited to L2 learner NWP and cloze (Section~\ref{sec:method}).
  \item A pilot experimental design on EFCAMDAT with zero-shot transfer
    to three external L2 corpora --- andrew100k (multi-L1, overlaps
    with EFCAMDAT), CELVA-SP (predominantly French-L1), and KUPA-KEYS
    (mixed-L1, several L1s out of EFCAMDAT's inventory) --- structured to isolate the gating
    contribution from data-only continual pre-training and from naive
    metadata concatenation (Section~\ref{sec:experiments}).
  \item A catalogue of nine complementary experiment variants
    (\texttt{experiments-ideas/v01--v09}) ablating the gate site
    ($G_1$ vs $G_2$), granularity (elementwise vs headwise), mode
    (multiplicative vs additive), metadata encoding (CEFR, L1, joint,
    typological), task (NWP vs cloze), and training regime (full vs
    frozen backbone).
\end{enumerate}

The claim we will argue is not that metadata-aware gating is a new
paradigm. It is the narrower one: the Qwen team's empirical finding ---
that the single most useful place to gate in softmax attention is at the
SDPA output --- happens to also be the place where learner-metadata
conditioning extracts the most leverage, because L2 effects are
fundamentally about \emph{which contextual information survives}
softmax-weighted aggregation. This conjecture has a testable prediction:
gate at $G_1$ beats the same gate at the value layer $G_2$, which in
turn beats gating at query or key projections --- the ordering Qwen
reports on native LM benchmarks should reproduce on learner NWP, and
the margin should be larger because L2 NWP places more demand on
attention filtering than native NWP.
